%%%%

%%%   Самарский государственный технический университет

%%%   Кафедра прикладной математики и информатики

%%%

%%%   Преамбула для статьи в журнал "Вестник Самарского государственного

%%%   технического университета. Серия: «Физико-математические науки»"

%%%   Ориентирована под MikTEx 2.4 for Win32

%%%

%%%   Хотя сам журнал собирается под GNU/Linux на дистрибутиве TeXLive



\documentclass[11pt,a4paper,twoside]{article}



\usepackage[cp1251]{inputenc}

\usepackage[T2A]{fontenc}

\usepackage[english,russian]{babel}

\usepackage{samgtubib}

\usepackage[dvips]{graphicx}

\usepackage[multiple]{footmisc}



\usepackage{samgtu}

\renewcommand{\VestnikYear}{2009} % Год издания Вестника

\renewcommand{\VestnikNumber}{2\,(19)} % Номер Вестника

\renewcommand{\VestnikMonth}{Ноябрь} % Месяц выхода Вестника

\renewcommand{\VestnikData}{01.11.09} % Дата подписания Вестника в печать

\renewcommand{\VestnikUPL}{26,64} % Усл. печ. л.

\renewcommand{\VestnikUIL}{25,73} % Уч.-изд. л.

\renewcommand{\VestnikRegNumber}{479/09} % Регистрационный номер

\renewcommand{\VestnikOrder}{994} % Номер заказа






%
%\begin{document}


\renewcommand{\baselinestretch}{0.88}

\renewcommand{\firstpage}{208}
\renewcommand{\lastpage}{211}

%\ShortPartition{Информатика}{}

\udk{519.178}\msc{68R10; 05C76}

\title{Структурный алгоритм распознавания предфрактального графа}
{Структурный алгоритм распознавания предфрактального графа}

\author{И.\,Х.~Утакаева}{У\,т\,а\,к\,а\,е\,в\,а~И.\,Х.}

\institute{Северо-Кавказская государственная гуманитарно-технологическая академия, \\
369000, Черкесск, ул.~Ставропольская, 36.}

\email{E-mail}{utakaev@yandex.ru}

\otherinf{\fullauthor{Ирина Хайрлыевна Утакаева}  \position{аспирант} \dept{каф. математики}}

\keywords{задача распознавания, предфрактальный граф}

\workabstract{Рассматривается задача распознавания предфрактального графа, образованного двумя полными
чередующимися затравками. Произведена математическая постановка, получены метрические характеристики, разработан эффективный
алгоритм распознавания, исследован вопрос о вычислительной сложности.}

\engtitle{Structural recognition algorithm for prefractal graph}

\engauthor{I.\,H.~Utakaeva}{U\,t\,a\,k\,a\,e\,v\,a~I.\,H.}

\otherenginf{\fullauthor{Irina H. Utakaeva}  \position{Postgraduate Student} \dept{Dept. of
Mathematics}}

\enginstitute{North-Caucasian State Humanitarian-Technological Academy, \\
36,  Stavropol'skaya, Cherkessk,  369000, Russia.  }

\engabstract{In the article the recognition problem for prefractal graph formed by two full alternating fuses is considered. Mathematical statement is made, metric characteristics are received, the effective algorithm of recognition is developed, the question on computing complexity is investigated.
}

\engkeywords{recognition problem, prefractal graph}

\date{17/III/2010}
\revisiondate{17/V/2011}

\maketitle

\small

%Одним из самых интересных свойств человеческого мозга является способность отвечать на бесконечное множество состояний внешней
%среды конечным числом реакций. Может быть, именно это свойство позволило человеку достигнуть высшей формы существования живой
%материи, выражающейся в способности к мышлению, т.\,е. активному отражению объективного мира в виде образов, понятий. Поэтому
%проблема распознавания возникла при изучении физиологических свойств мозга.

Проблема распознавания образов интересна как с прикладной, так и с принципиальной точки зрения. С прикладной точки зрения решение
этой проблемы важно прежде всего потому, что оно открывает возможность автоматизировать многие процессы, которые до сих пор
связывали лишь с деятельностью живого мозга. Принципиальное значение проблемы тесно связано с вопросом, который возникает с
развитием идей кибернетики: что может и что принципиально не может делать машина? В какой мере возможности машины могут быть
приближены к возможностям живого мозга? В частности, может ли машина развить в себе способность перенять у человека умение
производить определенные действия в зависимости от ситуаций, возникающих в окружающей среде? Пока стало ясно одно: если
человек может сначала сам осознать свое умение, то потом он его может описать. Если же человек обладает умением, но не может
объяснить его, то остается только один путь передачи умения машине "--- обучение примерами.

Круг задач, которые могут решаться с помощью распознающих систем, чрезвычайно широк. Сюда относятся не только задачи
распознавания зрительных и слуховых образов, но и задачи распознавания сложных процессов и явлений, возникающих, например, при
выборе целесообразных действий руководителем предприятия или выборе оптимального управления технологическими, экономическими,
транспортными или военными операциями. В каждой из таких задач анализируются некоторые явления, процессы, состояния внешнего
мира, всюду далее называемые объектами наблюдения. Прежде чем начать анализ какого"=либо объекта, нужно получить о нём
определенную, каким"=либо способом упорядоченную информацию. Такая информация представляет собой характеристику объектов, их
отображение на множестве воспринимающих органов распознающей системы \cite{utakaeva:bib1}.

Математической моделью многих задач распознавания является задача распознавания предфрактального графа \cite{utakaeva:bib2}.

Пусть представлен в явном виде некоторый граф  $G=(V,\,E)$, обладающий признаками предфрактального графа. Задача распознавания
предфрактального графа заключается в ответе на вопросы:
\begin{enumerate}
\item является ли данный граф предфрактальным с определенной затравкой \cite{utakaeva:bib2}?

\item можно ли построить эффективный алгоритм, который гарантированно построит процесс порождения предфрактального графа с
определенной затравкой?
\end{enumerate}


В дальнейшем будем использовать некоторые необходимые признаки предфрактальности графа $G=(V,\,E)$:
\begin{itemize}
\item[a)] для мощности множества вершин $| V |=N$ существует непустое множество пар $n$, $L$ таких, что $N=N(n,\,L)$;

\item[b)] для мощности множества рёбер $| E |$ существует хотя бы одна пара $n$, $L$, удовлетворяющая равенству $|Е_L|=q(n,\,L)$;

\item[c)] множество рёбер ранга $L$ состоит из объединения множеств рёбер затравок, появившихся в результате того, что каждая вершина
ранга $L-1$ графа была замещена затравкой.


\end{itemize}

В данной работе исследуется следующая задача.

\smallskip

\begin{task}Пусть задан в явном виде некоторый граф $G=(V,\,E),$ обладающий следующими необходимыми признаками предфрактального графа$:$
\begin{itemize}

\item[$a)$] для мощности множества вершин $| V |=N$ существует пара $n,$ $L$ такая$,$ что
$$
| V_L|=\begin{cases}
n^{\frac{L}{2}}(n+1)^{\frac{L}{2}}, & \text{при $L$ чётном}, \\
n^{\frac{L+1}{2}}(n+1)^{\frac{L-1}{2}},& \text{при $L$ нечётном};
\end{cases}
$$

\item[$b)$] для мощности множества рёбер $| E |$ существует пара $n,$ $L,$ удовлетворяющая равенству
$$
|E_L| = \begin{cases}
\left( q_1 + q_2 n \right) \frac{n^{\frac{L}{2}}(n+1)^{\frac{L}{2}}-1}{n(n+1)-1}, & \text{при $L$ чётном}, \\
\left( q_1 + q_2 n \right)
\frac{n^{\frac{L-1}{2}}(n+1)^{\frac{L-1}{2}}-1}{n(n+1)-1}+n^{\frac{L-1}{2}}(n+1)^{\frac{L-1}{2}}q_1,&\text{при $L$ нечётном},
\end{cases}
$$
где $q_1 =n(n+1)/{2},$ $q_2 = (n+1)(n+2)/{2};$

\item[$c)$] множество вершин состоит из двух подмножеств $V_1$ и $V_2,$ где $V_1 (V_2)$ "--- множество вершин $v\in V$ степени $\deg v=n+1$ $($степени $\deg v=n).$
\end{itemize}
\end{task}


Здесь излагается ответ на поставленный вопрос: является ли представленный граф  $G=(V,\,E)$ предфрактальным графом с
непересекающимися <<старыми>> ребрами, образованным двумя полными затравками $H_1 =(W_1,\,Q_1 )$ и $H_2 =(W_2,\,Q_2)$, где
мощности множеств вершин $|W_1|=n$, $|W_2|=n+1$? Процедура замещения вершины затравкой (ЗВЗ) \cite{utakaeva:bib2} производится на
нечётных номерах этапов $L$ затравкой $H_1 =(W_1,\,Q_1)$ и затравкой $H_2 =(W_2,\,Q_2)$ на чётных номерах этапов.

Для распознавания предфрактального графа $G=(V,\,E)$ построен алгоритм $\alpha _2$.

Описанию алгоритма предпошлём лемму.

\smallskip

\begin{lemma}
Пусть в предфрактальном графе $G=(V,\,E)$ две вершины $\nu_1$ и $\nu_2$ принадлежат одной затравке $H=(W,\,Q)$ $(\nu_1,\,\nu_2
\in W)$ и имеют смежность с некоторой вершиной $\nu \in V$, тогда $\nu$ также принадлежит этой затравке $H$.
\end{lemma}

\smallskip

Доказательство леммы осуществляется рассуждением от противного. Действительно, существование в предфрактальном графе $G=(V,\,E)$
двух вершин $\nu_1$ и $\nu_2$, принадлежащих одной затравке $H=(W,\,Q)$ $(\nu_1,\,\nu_2 \in W)$ и имеющих смежность с некоторой
вершиной $\nu \in V$ не своей затравки, означает, что в траектории предфрактального графа $G=(V,\,E)$ некоторый граф
$G_l=(V_l,\,E_l)$ содержит кратные ребра, т.\,е. является мультиграфом, что противоречит определению графа. Тогда вершина $\nu
\in V$ также принадлежит этой затравке $H$.

\smallskip

\Section[n]{Алгоритм $\boldsymbol{\alpha_2}$} Процедуру выделения затравки $H_1 =(W_1,\,Q_1)$ $(H_2 =(W_2,\,Q_2))$ обозначим через $\beta_1
(\beta_2)$. Исходя из принципа построения предфрактального графа будем применять ту или иную процедуру: если длина траектории
заданного предфрактального графа $L$ "--- чётная, то на последнем шаге было замещение затравкой $H_1 =(W_1,\,Q_1)$, поэтому на
первом этапе следует воспользоваться процедурой $\beta_1$, в противном случае "--- процедурой $\beta_2$. На последующих этапах
процедуры будут чередоваться.

\smallskip

\Section[n]{Описание процедуры $\boldsymbol{\beta _1}$ (выделение затравки $\boldsymbol{H_1 =(W_1,\,Q_1)}$)} Выделяем во множестве $V$ очередную
неотмеченную вершину $v_1 \in V$ \cite{utakaeva:bib2}. Если $v_1^\ast \in V_2$, то $\deg v_1^\ast =n$, т.\,е. $v_1^\ast$ смежна с
$(n-1)$ новыми вершинами своей затравки. Рассмотрим вершины, смежные с выделенной $v_k^\ast$, $k=2,\,3,\,\ldots,\,n+1$, но
неокрашенные, и выделим среди них те, которые будут иметь смежность между собой. По лемме 1 эти вершины также будут принадлежать
затравке. Если $v_1^\ast \in V_1$, то $\deg v_1^\ast =n+1$, т.\,е. $v_1^\ast$ смежна с $( n-1)$ новыми вершинами своей
затравки. Рассмотрим теперь вершины, смежные с выделенной $v_k^\ast$, $k=2,\,3,\,\ldots,\,n+2$, но неокрашенные, и выделим среди
них те, которые будут иметь смежность между собой. Согласно лемме~1, эти вершины также будет принадлежать затравке. Через $W'$
обозначаем множество всех вершин, отмеченных процедурой $\beta_1$. Если мощность $|W'|=n$, то выделяем и окрашиваем все ребра, у
каждого из которых концы представляют собой вершины данного множества $W'$. Работа процедуры $\beta_1$ завершается проверкой:
образует ли множество выделенных таким образом вершин и рёбер $n$-вершинный полный граф? Если да, то шаг, включающий в себя
описанную процедуру $\beta_1$, завершается результативно и следует переход к следующему шагу первого этапа. В противном случае
шаг считается безрезультатным и алгоритм $\alpha_2$ прекращает свою работу.

\smallskip

\Section[n]{Описание процедуры $\boldsymbol{\beta_2}$ (выделение затравки $\boldsymbol{H_2 =(W_2,\,Q_2)}$)} Выделяем во множестве $V$ очередную
неотмеченную вершину $v_1 \in V$. Если $v_1^\ast \in V_2$, то $\deg v_1^\ast =n$, т.\,е. $v_1^\ast $ смежна с $n$
новыми вершинами своей затравки, которые обозначаем через $v_k^\ast$, где $k=2,\,3\, \ldots,\,n+1$, и окрашиваем. Если же
$v_1^\ast \in V_1$, то $\deg v_1^\ast =n+1$, т.\,е. $v_1^\ast$ смежна с $n$ новыми вершинами своей затравки. Рассмотрим те
вершины, которые смежны с выделенной $v_k^\ast$, $k=2,\,4,\,\ldots,\,n+2$, но не окрашены, и выделим среди них те, которые будут
иметь смежность между собой. Согласно лемме~1, эти вершины также будет принадлежать этой затравке. Через $W'$ обозначаем
множество всех вершин, отмеченных процедурой $\beta_2$. Если мощность $| W'|=n+1$, то выделяем и окрашиваем все ребра, у каждого
из которых концы представляют собой вершины данного множества $W'$. Работа процедуры $\beta_2$ завершается проверкой: образует ли
множество выделенных таким образом вершин и рёбер $(n+1)$-вершинный полный граф? Если да, то шаг, включающий в себя описанную
процедуру $\beta_2$, завершается результативно и следует переход к следующему шагу первого этапа. В противном случае шаг
считается безрезультатным, и алгоритм $\alpha_2$ прекращает свою работу.

Этап $\rho =1$ начинает свою работу с проверки выполнения равенства
$$
|V_L|= \begin{cases}
n^{\frac{L}{2}}(n+1)^{\frac{L}{2}}, & \text{при $L$ чётном}, \\
n^{\frac{L+1}{2}}(n+1)^{\frac{L-1}{2}},& \text{при $L$ нечётном}.
\end{cases}
$$
Если равенство не выполняется, то алгоритм $\alpha_2$ заканчивает работу безрезультатно. В противном случае в графе $G$
выделяются множества $V_1$, состоящее из вершин степени $n+1$, и $V_2$, состоящее из вершин степени $n$. Если разность $V
\backslash ( V_1 \cup V_2) \ne \theta$, то алгоритм заканчивает работу с отрицательным результатом в том смысле, что данный граф
не является предфрактальным графом с непересекающимися <<старыми>> ребрами, образованным двумя полными затравками $H_1
=(W_1,\,Q_1)$ и $H_2 =(W_2,\,Q_2)$, где мощности вершин $|W_1|=n$, $|W_2|=n+1$, в котором процедура замещения вершины затравкой
(ЗВЗ) \cite{utakaeva:bib2} производится на нечётных этапах затравкой $H_1 =(W_1,\,Q_1)$, а на чётных "--- $H_2 =(W_2,\,Q_2)$. В
противном случае $V_1$ и $V_2$ образуют разбиение множества $V$, и дальнейшая работа этапа $\rho =1$ состоит из $m_0$ шагов, где
$m_0$ "--- число таких затравок, каждая из которых состоит из новых рёбер. Результатом каждого такого шага является выделенная в
графе $G$ очередная затравка.

В случае результативной работы каждого из $L-1$ этапов в качестве последнего члена последовательности получим $n$-вершинный
полный граф. Этот результат означает получение положительного ответа на вопрос: является ли представленный граф предфрактальным
графом с непересекающимися <<старыми>> ребрами, образованным двумя полными затравками $H_1 =(W_1,\,Q_1)$ и $H_2 =(W_2,\,Q_2)$,
где мощности вершин $|W_1|=n$, $|W_2|=n+1$, где процедура замещения вершины затравкой (ЗВЗ) производится на нечётных этапах
затравкой $H_1 =(W_1,\,Q_1)$, а на чётных $H_2 =(W_2,\,Q_2)$?

Распознаваемость исследуемого предфрактального графа $G=(V,\,E)$ вытекает из конструктивного описания алгоритма. Если полученную
последовательность $G_L^\ast$, $G_{L-1}^\ast$, $\ldots$,\, $G_1^\ast$ записать в обратном порядке, то имеет ли место совпадение
этой записи с траекторией? Ответ на этот вопрос можно считать положительным в том случае, если при построении этой
последовательности никогда не возникала альтернативность при выделении каждой затравки.

Рассмотрим вопрос о вычислительной сложности алгоритма $\alpha_2$. В процессе реализации этапов алгоритма $\alpha_2$
осуществляются следующие элементарные операции: определение степени вершины, выявление окрестности радиуса 1 для этой вершины,
просмотр всех пар вершин графа на предмет смежности, выделение и окрашивание рёбер. Поскольку все операции выполняются в пределах
одной затравки, верхняя оценка этапа не превосходит совокупного количества рёбер, выделенных и отмеченных в процессе работы
этапа. Отсюда справедлива следующая теорема.

\begin{theorem}Всякий предфрактальный граф $G_L =(V_L,\,E_L),$ образованный двумя полными чередующимися
затравками $H_1 =(W_1,\,Q_1 )$ и $H_2 =(W_2,\,Q_2),$ распознается алгоритмом $\alpha_2,$ если <<старые>> ребра не пересекаются, с
вычислительной трудоемкостью алгоритма $\tau (\alpha_2) \le O ( |E| L).$
\end{theorem}

\begin{thebibliography}{9}

\RBibitem{utakaeva:bib1}
\by Горелик~А.\,Л., Скрипкин~В.\,А.
\book Методы распознавания
\publaddr М.
\publ Высш. шк.
\yr 2004
\totalpages 261
\misctransl
\by Gorelik~A.\,L., Skripkin~V.\,A.
\book Pattern recognition methods
\publaddr Moscow
\publ Vyssh. shk.
\yr 2004
\totalpages 261


\RBibitem{utakaeva:bib2}
\by Кочкаров~А.\,М.
\book Распознавание фрактальных графов. Алгоритмический подход
\publaddr Нижний Архыз
\publ CYGNUS
\yr 1998
\totalpages 170
\misctransl
\by Kochkarov~A.\,M.
\book Recognition of fractal graphs. Algorithmic Approach
\publaddr Mizhniy Archyz
\publ CYGNUS
\yr 1998
\totalpages 170



\end{thebibliography}


\noindent
\hskip6.5cm \thedate \par\noindent
\hskip6.5cm\therevdate

\vspace{-10mm}

\makeengtitle

\renewcommand{\baselinestretch}{0.9}

%\end{document}
